% Options for packages loaded elsewhere
\PassOptionsToPackage{unicode}{hyperref}
\PassOptionsToPackage{hyphens}{url}
\documentclass[
]{article}
\usepackage{xcolor}
\usepackage[margin=1in]{geometry}
\usepackage{amsmath,amssymb}
\setcounter{secnumdepth}{-\maxdimen} % remove section numbering
\usepackage{iftex}
\ifPDFTeX
  \usepackage[T1]{fontenc}
  \usepackage[utf8]{inputenc}
  \usepackage{textcomp} % provide euro and other symbols
\else % if luatex or xetex
  \usepackage{unicode-math} % this also loads fontspec
  \defaultfontfeatures{Scale=MatchLowercase}
  \defaultfontfeatures[\rmfamily]{Ligatures=TeX,Scale=1}
\fi
\usepackage{lmodern}
\ifPDFTeX\else
  % xetex/luatex font selection
\fi
% Use upquote if available, for straight quotes in verbatim environments
\IfFileExists{upquote.sty}{\usepackage{upquote}}{}
\IfFileExists{microtype.sty}{% use microtype if available
  \usepackage[]{microtype}
  \UseMicrotypeSet[protrusion]{basicmath} % disable protrusion for tt fonts
}{}
\makeatletter
\@ifundefined{KOMAClassName}{% if non-KOMA class
  \IfFileExists{parskip.sty}{%
    \usepackage{parskip}
  }{% else
    \setlength{\parindent}{0pt}
    \setlength{\parskip}{6pt plus 2pt minus 1pt}}
}{% if KOMA class
  \KOMAoptions{parskip=half}}
\makeatother
\usepackage{color}
\usepackage{fancyvrb}
\newcommand{\VerbBar}{|}
\newcommand{\VERB}{\Verb[commandchars=\\\{\}]}
\DefineVerbatimEnvironment{Highlighting}{Verbatim}{commandchars=\\\{\}}
% Add ',fontsize=\small' for more characters per line
\usepackage{framed}
\definecolor{shadecolor}{RGB}{248,248,248}
\newenvironment{Shaded}{\begin{snugshade}}{\end{snugshade}}
\newcommand{\AlertTok}[1]{\textcolor[rgb]{0.94,0.16,0.16}{#1}}
\newcommand{\AnnotationTok}[1]{\textcolor[rgb]{0.56,0.35,0.01}{\textbf{\textit{#1}}}}
\newcommand{\AttributeTok}[1]{\textcolor[rgb]{0.13,0.29,0.53}{#1}}
\newcommand{\BaseNTok}[1]{\textcolor[rgb]{0.00,0.00,0.81}{#1}}
\newcommand{\BuiltInTok}[1]{#1}
\newcommand{\CharTok}[1]{\textcolor[rgb]{0.31,0.60,0.02}{#1}}
\newcommand{\CommentTok}[1]{\textcolor[rgb]{0.56,0.35,0.01}{\textit{#1}}}
\newcommand{\CommentVarTok}[1]{\textcolor[rgb]{0.56,0.35,0.01}{\textbf{\textit{#1}}}}
\newcommand{\ConstantTok}[1]{\textcolor[rgb]{0.56,0.35,0.01}{#1}}
\newcommand{\ControlFlowTok}[1]{\textcolor[rgb]{0.13,0.29,0.53}{\textbf{#1}}}
\newcommand{\DataTypeTok}[1]{\textcolor[rgb]{0.13,0.29,0.53}{#1}}
\newcommand{\DecValTok}[1]{\textcolor[rgb]{0.00,0.00,0.81}{#1}}
\newcommand{\DocumentationTok}[1]{\textcolor[rgb]{0.56,0.35,0.01}{\textbf{\textit{#1}}}}
\newcommand{\ErrorTok}[1]{\textcolor[rgb]{0.64,0.00,0.00}{\textbf{#1}}}
\newcommand{\ExtensionTok}[1]{#1}
\newcommand{\FloatTok}[1]{\textcolor[rgb]{0.00,0.00,0.81}{#1}}
\newcommand{\FunctionTok}[1]{\textcolor[rgb]{0.13,0.29,0.53}{\textbf{#1}}}
\newcommand{\ImportTok}[1]{#1}
\newcommand{\InformationTok}[1]{\textcolor[rgb]{0.56,0.35,0.01}{\textbf{\textit{#1}}}}
\newcommand{\KeywordTok}[1]{\textcolor[rgb]{0.13,0.29,0.53}{\textbf{#1}}}
\newcommand{\NormalTok}[1]{#1}
\newcommand{\OperatorTok}[1]{\textcolor[rgb]{0.81,0.36,0.00}{\textbf{#1}}}
\newcommand{\OtherTok}[1]{\textcolor[rgb]{0.56,0.35,0.01}{#1}}
\newcommand{\PreprocessorTok}[1]{\textcolor[rgb]{0.56,0.35,0.01}{\textit{#1}}}
\newcommand{\RegionMarkerTok}[1]{#1}
\newcommand{\SpecialCharTok}[1]{\textcolor[rgb]{0.81,0.36,0.00}{\textbf{#1}}}
\newcommand{\SpecialStringTok}[1]{\textcolor[rgb]{0.31,0.60,0.02}{#1}}
\newcommand{\StringTok}[1]{\textcolor[rgb]{0.31,0.60,0.02}{#1}}
\newcommand{\VariableTok}[1]{\textcolor[rgb]{0.00,0.00,0.00}{#1}}
\newcommand{\VerbatimStringTok}[1]{\textcolor[rgb]{0.31,0.60,0.02}{#1}}
\newcommand{\WarningTok}[1]{\textcolor[rgb]{0.56,0.35,0.01}{\textbf{\textit{#1}}}}
\usepackage{longtable,booktabs,array}
\newcounter{none} % for unnumbered tables
\usepackage{calc} % for calculating minipage widths
% Correct order of tables after \paragraph or \subparagraph
\usepackage{etoolbox}
\makeatletter
\patchcmd\longtable{\par}{\if@noskipsec\mbox{}\fi\par}{}{}
\makeatother
% Allow footnotes in longtable head/foot
\IfFileExists{footnotehyper.sty}{\usepackage{footnotehyper}}{\usepackage{footnote}}
\makesavenoteenv{longtable}
\usepackage{graphicx}
\makeatletter
\newsavebox\pandoc@box
\newcommand*\pandocbounded[1]{% scales image to fit in text height/width
  \sbox\pandoc@box{#1}%
  \Gscale@div\@tempa{\textheight}{\dimexpr\ht\pandoc@box+\dp\pandoc@box\relax}%
  \Gscale@div\@tempb{\linewidth}{\wd\pandoc@box}%
  \ifdim\@tempb\p@<\@tempa\p@\let\@tempa\@tempb\fi% select the smaller of both
  \ifdim\@tempa\p@<\p@\scalebox{\@tempa}{\usebox\pandoc@box}%
  \else\usebox{\pandoc@box}%
  \fi%
}
% Set default figure placement to htbp
\def\fps@figure{htbp}
\makeatother
\setlength{\emergencystretch}{3em} % prevent overfull lines
\providecommand{\tightlist}{%
  \setlength{\itemsep}{0pt}\setlength{\parskip}{0pt}}
\usepackage{bookmark}
\IfFileExists{xurl.sty}{\usepackage{xurl}}{} % add URL line breaks if available
\urlstyle{same}
\hypersetup{
  hidelinks,
  pdfcreator={LaTeX via pandoc}}

\author{}
\date{\vspace{-2.5em}}

\begin{document}

\section{Measurement Invariance Power
App}\label{measurement-invariance-power-app}

Shiny app and R functions for computing analytic power of chi-square
difference tests of \textbf{scale-level} (parameter-total) vs
\textbf{item-level} measurement invariance in multiple-group one-factor
CFA, generalizing the code behind the manuscript in
\texttt{Old\ Jordan\ Files\ for\ Claude/}.

\subsection{Run the app}\label{run-the-app}

\begin{Shaded}
\begin{Highlighting}[]
\NormalTok{shiny}\SpecialCharTok{::}\FunctionTok{runApp}\NormalTok{(}\StringTok{"path/to/ShinyApp"}\NormalTok{)}
\end{Highlighting}
\end{Shaded}

Requires: \texttt{lavaan}, \texttt{shiny}, \texttt{bslib},
\texttt{ggplot2}, \texttt{DT}.

\subsection{Files}\label{files}

{\def\LTcaptype{none} % do not increment counter
\begin{longtable}[]{@{}
  >{\raggedright\arraybackslash}p{(\linewidth - 2\tabcolsep) * \real{0.5000}}
  >{\raggedright\arraybackslash}p{(\linewidth - 2\tabcolsep) * \real{0.5000}}@{}}
\toprule\noalign{}
\begin{minipage}[b]{\linewidth}\raggedright
File
\end{minipage} & \begin{minipage}[b]{\linewidth}\raggedright
Contents
\end{minipage} \\
\midrule\noalign{}
\endhead
\bottomrule\noalign{}
\endlastfoot
\texttt{app.R} & Shiny UI/server only (no statistics live here) \\
\texttt{R/models.R} & Population specification
(\texttt{default\_params}, \texttt{apply\_noninvariance},
\texttt{build\_structured\_params}, \texttt{structured\_sweep},
\texttt{pop\_moments}) and lavaan model generation for any number of
items \emph{p} and groups \emph{G} (\texttt{make\_item\_model},
\texttt{make\_scale\_model}) \\
\texttt{R/power.R} & \texttt{mi\_power()} / \texttt{mi\_power\_sweep()}:
MacCallum-Browne-Cai (2006) analytic power for metric/scalar/strict
tests, both approaches, returned as tidy data frames \\
\texttt{R/plots.R} & \texttt{plot\_power()}: ggplot2 power curves (solid
= scale-level, dashed = item-level; color = overlaid condition) \\
\texttt{validation/validate\_manuscript.R} & Reproduces every population
RMSEA annotated in manuscript Figures 1-6 and 8 and checks power-formula
equivalence with the original pipeline. Run:
\texttt{Rscript\ validation/validate\_manuscript.R} (all 44 checks
pass) \\
\end{longtable}
}

\subsection{Scripting the functions
directly}\label{scripting-the-functions-directly}

\begin{Shaded}
\begin{Highlighting}[]
\FunctionTok{invisible}\NormalTok{(}\FunctionTok{lapply}\NormalTok{(}\FunctionTok{list.files}\NormalTok{(}\StringTok{"R"}\NormalTok{, }\AttributeTok{full.names =} \ConstantTok{TRUE}\NormalTok{), source))}
\FunctionTok{library}\NormalTok{(lavaan)}

\CommentTok{\# Manuscript Fig 3 (all panels in one plot): k = 7..2 non{-}invariant intercepts}
\NormalTok{conds }\OtherTok{\textless{}{-}} \FunctionTok{structured\_sweep}\NormalTok{(}
  \FunctionTok{list}\NormalTok{(}\AttributeTok{p =} \DecValTok{8}\NormalTok{, }\AttributeTok{n\_groups =} \DecValTok{2}\NormalTok{, }\AttributeTok{type =} \StringTok{"intercepts"}\NormalTok{, }\AttributeTok{n\_ni =} \DecValTok{3}\NormalTok{, }\AttributeTok{magnitude =}\NormalTok{ .}\DecValTok{1}\NormalTok{),}
  \AttributeTok{sweep\_var =} \StringTok{"n\_ni"}\NormalTok{, }\AttributeTok{values =} \DecValTok{7}\SpecialCharTok{:}\DecValTok{2}\NormalTok{)}
\NormalTok{res }\OtherTok{\textless{}{-}} \FunctionTok{mi\_power\_sweep}\NormalTok{(conds, }\AttributeTok{tests =} \FunctionTok{c}\NormalTok{(}\StringTok{"metric"}\NormalTok{, }\StringTok{"scalar"}\NormalTok{))}
\FunctionTok{plot\_power}\NormalTok{(res}\SpecialCharTok{$}\NormalTok{curves)}
\NormalTok{res}\SpecialCharTok{$}\NormalTok{details   }\CommentTok{\# population RMSEA, test df, N needed for 80\% power}
\end{Highlighting}
\end{Shaded}

\subsection{Reproducing the manuscript figures in the
app}\label{reproducing-the-manuscript-figures-in-the-app}

Figure numbers refer to ``figures with captions.pdf'' (the output of
\texttt{figure\ markdown4.Rmd}). Each 6-panel manuscript figure becomes
\textbf{one} plot in the app: the compared values (colored curves, set
in the ``Separate curves'' panel) correspond to the panels (a)-(f). Read
the facet named in the ``Read'' column; the other facet shows the same
population under the other test.

Common settings for every figure (the app defaults): \emph{Model}: 2
groups, p = 8 items, base loading .7, base residual variance .51.
\emph{Additional settings}: N per group 50-1000, alpha = .05. Tick
``Also test strict invariance'' for the residual-variance figures (it
turns on automatically when the non-invariance type is set to residual
variances).

The manuscript places non-invariance on different item positions than
the app (which always uses the last k items); because all other items
are exchangeable this yields identical fit statistics and power.

{\def\LTcaptype{none} % do not increment counter
\begin{longtable}[]{@{}
  >{\raggedright\arraybackslash}p{(\linewidth - 6\tabcolsep) * \real{0.2500}}
  >{\raggedright\arraybackslash}p{(\linewidth - 6\tabcolsep) * \real{0.2500}}
  >{\raggedright\arraybackslash}p{(\linewidth - 6\tabcolsep) * \real{0.2500}}
  >{\raggedright\arraybackslash}p{(\linewidth - 6\tabcolsep) * \real{0.2500}}@{}}
\toprule\noalign{}
\begin{minipage}[b]{\linewidth}\raggedright
Fig
\end{minipage} & \begin{minipage}[b]{\linewidth}\raggedright
Non-invariance panel
\end{minipage} & \begin{minipage}[b]{\linewidth}\raggedright
Separate-curves panel
\end{minipage} & \begin{minipage}[b]{\linewidth}\raggedright
Read
\end{minipage} \\
\midrule\noalign{}
\endhead
\bottomrule\noalign{}
\endlastfoot
1 (a,b) & Type: Loadings; count: 6; magnitude: per-item .1; direction:
\textbf{Alternating +/-} & One curve per value of: Magnitude; values:
\texttt{0.1,\ 0.2} & Metric facet: item-level (dashed) has power,
scale-level (solid) stays at .05 \\
1 (c,d) & as above but Type: Intercepts & same & Scalar facet \\
1 (e,f) & as above but Type: Residual variances & same & Strict facet \\
2 & Type: Loadings; magnitude: per-item .1; direction: All positive &
One curve per value of: Number of non-invariant items; range 2-7 &
Metric facet (panels a-f = 7,6,5,4,3,2 non-invariant) \\
3 & Type: Intercepts; otherwise as Fig 2 & same & Scalar facet \\
4 & Type: Residual variances; otherwise as Fig 2 & same & Strict
facet \\
5 & Type: Loadings; magnitude: \textbf{Total across non-invariant items}
= .6; direction: All positive & One curve per value of: Number of
non-invariant items; range 1-6 & Metric facet (per-item difference
becomes .6/k: .1, .12, .15, .2, .3, .6) \\
6 & Type: Intercepts; otherwise as Fig 5 & same & Scalar facet \\
7 & Type: Residual variances; otherwise as Fig 5 & same & Strict
facet \\
8 & Type: Loadings; count: 3; magnitude: per-item .2; direction: All
positive & One curve per value of: Number of items (p); values:
\texttt{4,\ 6,\ 8,\ 10,\ 12,\ 14} & Metric facet \\
9 & Type: Intercepts; otherwise as Fig 8 & same & Scalar facet \\
10 & Type: Residual variances; otherwise as Fig 8 & same & Strict
facet \\
\end{longtable}
}

Notes: - Fig 1's panels differ only in magnitude (.1 vs .2), so the app
shows the two panels of each pair as two colored curves in one facet.
With the alternating direction and an even count the parameter totals
are equal across groups, so composite comparisons are unaffected, which
is the point of Study 1. - The manuscript curves show power of the test
matching the manipulated parameter type only; the app additionally plots
the other tests for the same population. For Fig 2/5/8 (loadings) the
scalar facet is flat at alpha because intercepts remain invariant, and
vice versa. - Exact expected rmsea values for Figs 1-6 and 8 are
asserted in \texttt{validation/validate\_manuscript.R} (the Details
tab's ``RMSEA (model)'' column should match the figure annotations to 3
decimals).

\subsection{Method notes / differences from the original
scripts}\label{method-notes-differences-from-the-original-scripts}

\begin{itemize}
\tightlist
\item
  Noncentrality is computed from the population minimum discrepancy:
  \texttt{nonc\ =\ (N\_total\ -\ G)\ *\ (F0\_constrained\ -\ F0\_baseline)},
  which for G = 2 is numerically identical to the manuscript's
  \texttt{n\ *\ df\ *\ rmsea\^{}2} but is correct for any G and
  independent of lavaan's multi-group RMSEA convention. Baselines are
  one invariance level lower (configural -\textgreater{} metric
  -\textgreater{} scalar -\textgreater{} strict), fit with the same
  approach.
\item
  Test df are taken from the fitted models. This silently fixes a bug in
  the original \texttt{power.plot.doer.weak}/\texttt{.strict}, where the
  critical value used \texttt{nrow(pop.cov1-1)} = \emph{p} df instead of
  \emph{p - 1} for the metric test (a small effect, e.g.~power .797 vs
  .819 at n = 300 in Fig 1a).
\item
  Scale-level models for G groups impose G - 1 total-equality
  constraints per parameter class (plus the first-intercept anchor for
  scalar).
\item
  Population latent variance is 1 and latent mean 0 in all groups;
  ad-hoc effect sizes (beta\_spurious, d\_spurious) are intentionally
  omitted.
\end{itemize}

\end{document}
